\section{Introducción}

El problema del signo de la forma de Weil reúne dos exigencias:
conservar todas sus contribuciones aritméticas y controlar la forma
en espacios completos de funciones. El criterio global relaciona
su positividad con la hipótesis de Riemann. Una construcción
espectral que aspire a ese criterio debe, por tanto, especificar
el producto, los pesos y los dominios que determinan la forma,
y demostrar cómo se transmiten sus términos cruzados. Este
artículo desarrolla esa relación desde una aritmética de caminos
con memoria y obtiene resultados exactos de reconstrucción,
transporte y coercividad.

La contribución funcional principal es una reducción de Schur
para cada ventana fijada. Una partición suficientemente fina
separa las medias celulares de un espacio infinito de detalles
de media nula cuya forma es coerciva. La eliminación exacta de
estos detalles conserva todo su acoplamiento con las medias y
produce una matriz finita con el mismo índice negativo que la
forma completa. El resultado da un criterio exacto, no una
sustitución del espacio de funciones por muestras. Su aplicación
en la ventana $(-1/18,1/18)$ demuestra una cota inferior de
$\frac12\|f\|_2^2$ para todo el dominio conjunto de sus lectores.
Un transporte reversible prolonga esta cota a colecciones completas
de pruebas con colas exponenciales.

\subsection*{Construcción aritmética y estructura conservada}

El punto de partida es un núcleo formal matemático denominado
Holografía Modular Triádica, en adelante HMT. Su objeto
\emph{All Possible Paths} (APP) reúne caminos y evaluaciones
aritméticas en el soporte $X_9=(\mathbb Z/9\mathbb Z)^2$.
Los generadores cardinales $\{(1,0),(-1,0),(0,1),(0,-1)\}$
definen el grafo de Cayley de su grupo aditivo, una retícula
cíclica bidimensional con realización toroidal. Sobre las mismas
posiciones actúan las evaluaciones suma y producto, cada una
con su residuo y su cociente. El soporte especifica dónde se
desplaza una ruta; las hojas aritméticas especifican qué se
calcula y conserva durante ese desplazamiento.

TRIT es la firma ternaria orientada $\{-1,0,+1\}$, obtenida
mediante representantes balanceados módulo $3$. Distingue los
tipos hiperbólico, parabólico y elíptico en su realización
cuadrática. El \emph{Topological Phase Kernel}, en adelante
TPK, compone selección de fase, transporte y actualización de
memoria. Su estado conserva posiciones, hojas, orientación,
cocientes, prefijos y frontera. La prolongación nonádica transporta
esta información entre profundidades: el retorno de fase avanza
la historia en vez de reiniciarla.

La estructura discreta conjunta del continuo se desarrolla sobre
esos mismos estados y sus compatibilidades. Sus construcciones
correlativas conservan restricciones, incidencias, formas y
memoria dentro de un único sistema de refinamiento. La lectura
arquimediana obtiene valores a partir de cilindros compatibles y la lectura
de frontera retiene las incidencias utilizadas por las
realizaciones posteriores. El presente problema aritmético y
espectral conserva ese origen común, con las operaciones y los
dominios que intervienen efectivamente en cada resultado.

La relación entre centro y frontera proporciona una primera
instancia de reconstrucción. En la coordenatización multiafín de APP, la
frontera levantada determina el interior; los valores
residuales se completan con los cocientes que distinguen sus
valores enteros. Las firmas de continuación requieren conservar
las hojas, y el calendario de vacancias conserva la discrepancia
entre capacidad y profundidad. La
sección~\ref{sec:centro-frontera} sitúa en esa construcción el
centro $(5,5)$, la firma regional $555555$, los lectores
$72,27,77,22$ y los intervalos $21/22$ y $6/7$. Su función
especifica la información transportable antes de cualquier
identificación espectral.

\subsection*{Del producto a la función completada}

El producto nativo compone las dos hojas de APP mediante un
transductor de dígitos y cocientes. Su invariante de evaluación
prueba la exactitud de la operación sobre palabras nonádicas,
y la unicidad de la forma normal permite reconocer después la
multiplicación entera. Las fibras del producto determinan sus
aperturas; la ausencia de una apertura no trivial selecciona
los irreducibles. El teorema~\ref{rz-num-01-6} identifica sus
evaluaciones con los primos positivos.

La factorización y los relojes aritméticos describen propiedades
distintas de esos objetos. La primera determina sus componentes
multiplicativos; los segundos registran retornos de órbitas
residuales. La representación por ocupaciones de modos conserva
la multiplicidad de cada irreducible y conduce al producto de
Euler. La traza sobre todas las ocupaciones y la traza sobre
los modos individuales se relacionan mediante los momentos
primo--potencia de la proposición~\ref{rz-num-02-3}.

El refinamiento cíclico proporciona a continuación una traza de
calor cuyo límite es la serie theta. La transformación de Mellin
produce $\pi^{-s/2}\Gamma(s/2)Z(s)$; la inversión de la escala
determina la contribución polar y la ecuación funcional de la
forma completada. Las proposiciones~\ref{rz-num-03-1}
y~\ref{rz-num-03-2} mantienen la normalización y la convergencia
de ese paso. Los desarrollos de partes finitas y lectores
residuales amplían sus evaluaciones con restos controlados.
La forma de Weil conserva conjuntamente los términos primo,
gamma y polar que proceden de esta composición.

\subsection*{Doble círculo, reconstrucción y transporte}

El doble círculo distingue la coordenada espectral de la memoria
de vueltas. La descomposición de nueve fases produce un balance
exacto con coeficiente $8/81$; los modos del bloque central
determinan la razón $q=5/9$. Los momentos de memoria
$q^{|n|}$ poseen una factorización triangular positiva en toda
dimensión. El registro de pérdidas conserva además el término
terminal y los productos cruzados, de modo que la contracción
visible puede acompañarse de una conservación isométrica de la
historia registrada.

La realización analítica construye, sobre pruebas con
decaimiento exponencial, un operador reversible $C_a$, con
$a>b>1/2$. Su acción integral conserva simultáneamente las
correlaciones y los productos polares cruzados. La identidad
de la subsección~\ref{tm-add:invariancia-weil} demuestra
\[
\mathscr W(C_af,C_ag)=\mathscr W(f,g).
\]
El margen $a>1/2$ pertenece al dominio de este lector analítico
y garantiza la convergencia de sus contribuciones polares.
El círculo regular de memoria y una representación que recibe
un límite atómico conservan sus dominios respectivos.

La cola gamma contiene un lector inverso explícito: recupera
la función de prueba y permite reconstruir la columna negativa
desde la positiva. Las ecuaciones~\eqref{gamma-add-eq-12}
y~\eqref{gamma-add-eq-14} describen su acción, y
\eqref{gamma-add-eq-16} conserva la diferencia completa de
formas. Esta reconstrucción da el operador que debe estudiarse;
su acotación y su contractividad son propiedades diferentes,
cuyas estimaciones se desarrollan a continuación.

\subsection*{Coercividad y criterio de prolongación global}

El teorema~\ref{schur-add-schur-exacto} demuestra que, una vez
obtenida la cota positiva del bloque de detalles $D$, la positividad
de la forma en $V_R$ equivale a la de
$S=A-B^*D^{-1}B$. La corrección conserva el efecto de todos
los detalles y sus cruces con las medias. El
teorema~\ref{schur-add-residual} encierra esa matriz mediante
residuos y una cota operatoria del error. En la ventana de
longitud $1/9$, las estimaciones
$D>I$, $A>97/100$ y $\|B\|^2<1/5$ conducen al
teorema~\ref{schur-add:coercividad-ventana}. La ventana local
satisface $1/9<\log2$; sus colecciones transportadas pueden
tener colas y reciben entonces la suma prima completa mediante
la invariancia demostrada.

El criterio global exige también los términos cruzados entre
colecciones distintas y la cobertura de todo el dominio de prueba.
La formulación precisa se encuentra en las
subsecciones~\ref{tm-add-m6-2}, \ref{tm-add-m7}
y~\ref{gamma-add-sub-375}: identificar los momentos
aritméticos completos con un Gram positivo equivale a conservar
simultáneamente la energía y el transporte, con sus dominios
y límites. Esa condición se mantiene como hipótesis de los
resultados globales; la coercividad demostrada pertenece a los
dominios explícitos anteriores. La aportación del artículo es
unir la generación aritmética con identidades de reconstrucción
y estimaciones funcionales suficientemente precisas para
establecer ese alcance y su condición exacta de prolongación.
\clearpage
